\chapter*{Záver}  % chapter* je necislovana kapitola
\addcontentsline{toc}{chapter}{Záver} % rucne pridanie do obsahu
\markboth{Záver}{Záver} % vyriesenie hlaviciek

Podarilo sa nám vytvoriť nové nástroje obohacujúce aplikáciu Prieskumník štruktúr o nové možnosti vizualizácie a tvorby vzťahov.

V prvom rade sme sa venovali návrhu grafových editorov, ktoré sa zatiaľ ukázali byť prínosné pri sledovaní vzťahov v rámci interpretácií binárnych predikátov a unárnych funkcií. Snažili sme sa dosiahnuť prehľadnú a prispôsobiteľnú vizualizáciu. Jednotlivé reprezentácie sme navrhli tak, aby spĺňali všetky požadované vlastnosti. Tiež sme sa venovali ďalším vylepšeniam aplikácie. Maticový a Databázový editor ponúkajú ďalší alternatívny pohľad na interpretácie predikátov a funkcií. Editor interpretácií funkcií rozborom prípad umožňuje jednouchšie vytváranie interpretácií funkcií, čo poskytuje učiteľom predmetu možnosť vytvárať úlohy, kde by využitie predošlého riešenia bolo veľmi nepraktické. Tiež sme pridali novú sekciu aplikácie slúžiacu na definovanie dopytov na splniteľnosť formúl.

Po návrhu sme popísali implementáciu jednotlivé riešení do existujúcej aplikácie. Zabezpečili sme jej kompatibilitu s Logickým pracovným zošitom a nakoniec sme sa zamerali na vylepšenia týkajúce sa Henkinovej-Hintikkovej hry.

Aplikáciu sa tiež podarilo dostať pomerne skoro do stavu, kedy mohla byť nasadená do výučby. Žiaci mali počas semestra príležitosť jej rôznu funkcionalitu využiť pri riešení cvičení a domácich úloh. Ku koncu výučby sme využili jedno z cvičení na podrobnejšie otestovanie nových nástrojov, po ktorom sme pomocou dotazníka od študentov získavali spätnú väzbu.

Pri jednotlivých nástrojoch je však určite stále priestor na zlepšenie. Jedným z nich je pri zobrazení bipartitného grafu umožniť separátne filtrovanie zobrazenej domény oboch skupín. V momentálnom stave je doména zobrazená v oboch skupinách identická, čo často nie je dostatočne flexibilné na dosiahnutie uspokojivého zobrazenia. Ďalším potenciálnym zlepšením je pridať možnosť konfigurovať farebnú paletu používanú pri vizualizácií unárnych predikátov. Rozhranie editora interpretácií funkcií môže tiež novým používateľom prísť na prvý pohľad neintuitívne, čo by sa dalo vylepšiť zmenou jeho rozhrania, hlavne v oblasti pridávania novej podmienky.

%Podarilo sa nám jasne stanoviť ciele a požiadavky na nové grafové editory pre prieskumník štruktúr logiky prvého rádu. Jednotlivé grafové reprezentácie sa zatiaľ ukázali byť vhodné a dobre sa navzájom dopĺňajúce na dosiahnutie zmysluplných, intuitívnych a vysoko informatívnych vizualizačných prostriedkov pre študentov.
%
%Ďalej sme sa v prvom návrhu zaoberali nápadmi, ako premeniť tieto reprezentácie do interaktívnej podoby a ešte lepšie využiť ich potenciál s ohľadom na naše špecifické požiadavky. Nakoniec sa ukázala byť prínosná aj pilotná implementácia, ktorá v mnohých ohľadoch utvrdila správnosť myšlienok predstavených v predchádzajúcich kapitolách, no zároveň odhalila nedostatky niektorých riešení a potrebu ich vylepšenia v ďalších iteráciách.
%
%Pri mojom vlastnom testovaní sa prejavila najmä slabá intuitívnosť častých úkonov, ako napríklad pridávania alebo mazania nových vrcholov grafov. K tomuto problému tiež prispieva fakt, že na vzor prvotného návrhu je táto funkcionalita implementovaná pre bipartitný graf inak, ako pre ostatné dve reprezentácie. Tento problém by som chcel vyriešiť kompletnou zmenou a zjednotením spôsobu ich pridávania, napríklad implementáciou mechanizmu filtrovania vrcholov, podobnému tomu pri volení unárnych predikátov v orientovanom grafe. V niektorých aspektoch je zas zrejmé, že vizuálna podoba pilotnej implementácie sa dostatočne nedrží prvotného návrhu, na čom bude ešte tiež potrebné zapracovať.
%
%Medzi menej viditeľné, no o nič menej zásadné problémy patrí nutnosť doladiť spôsob, akým jednotlivé grafy spolu komunikujú zmeny týkajúce sa ich spoločného stavu. Pri tejto miere interaktivity, ktorú knižnica React umožňuje, je dôležité dbať aj na efektívnosť implementovaných riešení. V opačnom prípade to môže viesť k zlému zážitku pri interakcii s aplikáciou, z dôvodu neresponzívneho používateľského rozhrania. Tento aspekt bude predstavovať jeden z mojich hlavných cieľov.



